\documentclass[UTF8,11pt]{ctexart}
\usepackage[a4paper,margin=2.2cm]{geometry}
\usepackage{hyperref}
\usepackage{enumitem}
\usepackage{longtable}
\hypersetup{colorlinks=true,linkcolor=black,urlcolor=blue}
\setlength{\parskip}{0.5em}
\setlength{\parindent}{2em}
\title{《迈出理解闪电贷及其在 DeFi 生态中应用的第一步》\\中文译读版}
\author{原文作者：Dabao Wang, Siwei Wu, Ziling Lin, Lei Wu, Xingliang Yuan, Yajin Zhou, Haoyu Wang, Kui Ren}
\date{原文来源：SBC 2021\\DOI：\url{https://doi.org/10.1145/3457977.3460301}}
\begin{document}
\maketitle

\section*{说明}
本文档依据用户提供的 PDF《Towards A First Step to Understand Flash Loan and Its Applications in DeFi Ecosystem》整理为中文译读版，目标是便于课程项目阅读与引用。内容按原论文结构进行中文转写与重组，重点保留论文的研究问题、方法、核心结果、应用分类与结论；版式不保留原始双栏排版，个别公式图表改为中文说明。

\section*{摘要}
闪电贷是去中心化金融生态中出现的一类新型服务，它允许用户在不提供抵押物的前提下，在单笔区块链交易内部完成借款与还款。该机制一方面为套利、清算和资产置换提供了极高的资金效率，另一方面也使攻击者能够临时调动本不拥有的大额资金，从而放大恶意操作的影响。尽管媒体已经报道过多起利用闪电贷的攻击事件，但当时学术界对于现有闪电贷服务本身的工作方式、使用规模以及典型用途，仍缺少系统性的认识。

本文以三个主流平台提供的闪电贷服务为对象，首先说明闪电贷提供方与用户之间的交互流程，然后针对不同平台设计三类链上识别模式，以从历史交易中定位闪电贷交易。基于这些模式，作者从截至 2021 年 1 月 31 日的以太坊交易中识别出 76,303 笔闪电贷交易。论文进一步分析了闪电贷随时间增长的使用趋势，并总结出四类典型应用：套利、刷量交易、闪电清算和抵押品置换。最后，作者提出了两个值得继续研究的方向：一类面向套利自动化，另一类面向 DeFi 攻击检测。

\section{引言}
DeFi 在近几年快速增长，大量借贷、交易、清算和资产管理活动开始由智能合约完成。闪电贷在这种环境下成为一种极具代表性的基础能力：用户无需预存抵押物，只要能在同一笔交易结束前偿还本金和费用，就可以临时借到较大规模的资产。

论文指出，闪电贷并不是“天然恶意”的工具。对正常交易者而言，它可以显著降低套利和资产再配置的门槛。例如，当不同 DEX 上同一资产出现价差时，交易者不需要提前持有大额资金，只需借入资金、执行一系列交换、再在交易结束前归还借款，就能完成套利。对清算者而言，闪电贷也使“没有大量自有资金的人”可以参与 DeFi 借贷市场中的清算活动。

但与此同时，闪电贷也会显著放大系统脆弱性。如果某个协议的价格预言机、会计逻辑或交易路径设计存在缺陷，攻击者就可以在极短时间内调动大额资金操纵市场条件，从而把原本难以利用的问题变成可盈利的攻击。论文以 bZx 与 Eminence 等案例为背景，说明当时已经出现了因闪电贷而被放大的严重损失。

作者认为，为了真正讨论防御手段，首先必须系统回答三个问题：第一，主流闪电贷协议的工作流程是什么；第二，如何从以太坊海量历史交易中识别闪电贷交易；第三，这些交易究竟被用在什么场景中。因此，这篇论文的定位是一次“测量与理解”的起点工作。

\section{闪电贷交互流程}
论文将闪电贷的一般流程拆解为五步。第一步，用户部署或准备一个接收闪电贷的合约。第二步，该合约向闪电贷提供方发起借款请求。第三步，协议先把请求的资产转给用户合约，并立即回调用户定义的逻辑入口。第四步，用户在该回调函数内部执行套利、清算、兑换或其它自定义操作。第五步，在交易结束前，借入资产及相应费用必须返还给协议；如果返还失败，整笔交易回滚。

从这个角度看，闪电贷的本质并不是“无风险借钱”，而是“利用区块链交易的原子性，把短时间的流动性访问能力嵌入到一笔交易里”。原子性意味着：要么所有步骤都成功并完成还款，要么整个交易作废，因此协议无需像传统借贷那样要求抵押物。

\section{三类闪电贷提供方}
\subsection{Aave}
Aave 提供原生的闪电贷接口。用户需要编写接收者合约，并在其中实现回调逻辑。调用流程通常包括：向 Aave Pool 发起借款请求、Pool 将资产转出、Pool 回调用户合约执行操作、用户最终归还本金与手续费。如果协议检测到余额未恢复，就会回滚交易。论文把 Aave 视为最典型、最直接的闪电贷实现。

\subsection{dYdX}
论文研究时，dYdX 并没有单独暴露名为 flash loan 的接口，但它可以通过一组组合操作实现等价效果。用户可以在同一笔交易中调用 \texttt{operate}，并通过 \texttt{withdraw}、\texttt{callFunction}、\texttt{deposit} 等动作形成“借出-执行-归还”的闭环。由于 dYdX 当时不对这类调用收取专门的闪电贷费用，因此它在某些用例中具有成本优势。

\subsection{Uniswap V2}
Uniswap V2 的闪电贷能力不是单独设计的借贷接口，而是内嵌在 \texttt{swap} 机制中的 flash swap。用户调用 pair 合约的 \texttt{swap(amount0Out, amount1Out, to, data)} 时，如果 \texttt{data} 非空，pair 会先把 token 转出，再触发回调函数，由调用者在回调中执行逻辑并补足 pair 所需资产。与 Aave 和 dYdX 不同，Uniswap V2 的逻辑更像“先拿到资产，再在交易末尾恢复 pair 的经济不变量”。论文还指出，相较于其他提供方，Uniswap V2 的收费更高，约为借出资产的 0.3\%。

\section{闪电贷识别模式}
论文的核心贡献之一，是根据不同协议的实现特征设计识别模式，而不是试图用一个完全统一的黑盒标准去检测全部交易。

\subsection{模式一：Aave}
对 Aave 而言，识别的关键依据是其闪电贷执行成功后会产生专门的 \texttt{FlashLoan} 事件。只要一笔交易中出现该事件，并且事件中的关键字段满足论文定义的条件，就可以将其识别为 Aave 闪电贷交易。由于事件语义清晰，这一模式的识别难度相对较低。

\subsection{模式二：dYdX}
对 dYdX 而言，识别不能只依赖单一事件，而需要观察一组动作顺序。论文将 \texttt{operate} 中出现的 \texttt{withdraw}、用户回调与 \texttt{deposit} 归还过程视为一个典型模式。当相关操作在同一笔交易中形成完整闭环时，即可判定该交易使用了 dYdX 提供的闪电贷能力。

\subsection{模式三：Uniswap V2}
对 Uniswap V2 而言，识别重点在于 flash swap 的执行结构：pair 先转出 token，随后调用接收方回调，接收方完成操作后必须在同一笔交易中把所需资产返还给 pair。论文据此总结了识别要点，包括 \texttt{swap(..., data)} 的调用形式、内部转账以及资金回流 pair 的路径。

作者特别强调，这些模式强依赖协议特定实现，因此扩展到新协议时需要重新设计新的检测规则。这一结论对后续课程项目尤其重要：规则型检测器具备很强可解释性，但天然依赖协议知识。

\section{闪电贷交易测量结果}
作者从以太坊账本中收集了大约十亿笔交易，并将上述三类识别模式应用到截至 2021 年 1 月 31 日的历史数据。最终识别出 76,303 笔闪电贷交易，以及 1,454 个闪电贷接收者地址。

从协议分布看，Uniswap V2 的闪电贷交易数量最多，共 36,574 笔，占总量约 48\%；dYdX 为 24,983 笔，占约 32.5\%；Aave 为 15,016 笔，占约 19.5\%。从接收者数量看，dYdX 对应的独立接收者最多；从平均使用强度看，Uniswap V2 上每个接收者发起的闪电贷次数最高。论文据此说明，不同协议虽然都提供“单交易借款”能力，但它们在实际生态中的使用方式并不一致。

论文还给出了时间趋势。无论按绝对数量还是按“闪电贷交易数占总交易数的比例”衡量，闪电贷都呈持续增长趋势。作者据此得出结论：闪电贷正在成为 DeFi 中越来越常见的基础操作，不应仅仅被视作极少数攻击案例的附属现象。

\section{四类典型应用}
\subsection{套利}
套利是论文给出的第一类典型应用。不同平台上同一资产价格可能出现短时间偏差，交易者可以利用闪电贷借入资金，在多个池子之间快速买低卖高，再在交易结束前归还借款。论文指出，套利本身并不必然是恶意行为，它在某种程度上还能帮助市场价格更快收敛。作者展示了一笔真实交易：用户从 dYdX 借入 1.13 ETH，在 Balancer 与 Uniswap 之间完成资产交换，扣除 gas 后仍获得正收益。

\subsection{刷量交易}
第二类应用是 wash trading，即人为制造交易量。论文指出，这种行为会夸大某些池子或资产的活跃度，误导普通用户。闪电贷在这里的作用是降低作恶门槛：操作者不需要提前持有大量资产，只要能承担 gas 费和可能的小额亏损，就能通过反复买卖在短时间内制造大量成交量。论文举例说明，有用户从 Aave 借入 ETH，在一个 Uniswap 池子中多次来回兑换，以提升该池子表面交易量。

\subsection{闪电清算}
第三类应用是 liquidation。借贷协议中的头寸一旦抵押不足，就可能被清算者用折扣价格接管。传统上，这要求清算者手头已有足够资本；有了闪电贷之后，任何人只要能在一笔交易中完成借款、买入、处置和还款，就能参与清算并获取折扣收益。论文举例展示了一笔利用 dYdX 闪电贷在 Compound 中执行清算的真实交易，清算者在完成兑换与归还后获得了数百 DAI 的利润。

\subsection{抵押品置换}
第四类应用是 collateral swap。它通常包含两步：先赎回旧抵押物，再把赎回的抵押物转换为新的更稳定或更适合当前市场状态的资产。由于市场波动剧烈，用户有时需要快速关闭旧头寸、替换抵押物，以避免被清算或遭受价格滑点。闪电贷使用户无需提前持有额外资金，也能在一笔交易里完成赎回、转换与再归还。论文给出的例子中，用户借入 Aave 的 DAI，用来赎回原先抵押的 BAT，再把 BAT 转成更稳定的 USDC，最后用其中一部分归还闪电贷，从而在没有额外 DAI 的情况下完成抵押物调整。

\section{未来研究方向}
论文最后提出两个值得继续深入的研究方向。

第一，\textbf{套利自动化}。由于 DeFi 市场几乎每天都有价格偏差，越来越多组织和个人使用机器人监测价差并自动发起闪电贷套利。作者认为，如何更及时地检测可套利信息并把它输入自动化交易系统，是一个重要方向。

第二，\textbf{DeFi 攻击检测}。随着 DeFi 普及，攻击者可能通过操纵预言机、组合多个协议或利用设计缺陷窃取资金。尤其是零日攻击，由于涉及多个实体和复杂调用路径，传统方法难以及时识别。如何提出更有效的方法去检测恶意交易，仍然是一个开放问题。

\section{结论}
本文是对闪电贷生态的早期系统研究。它并不试图解决所有防御问题，而是首先回答“闪电贷如何工作、如何识别、被用于做什么”这三个基础问题。作者基于三类协议交互模式识别出 76,303 笔闪电贷交易和 1,454 个接收者，并通过真实案例归纳了套利、刷量、清算和抵押品置换四种主要用途。论文的贡献在于把闪电贷从零散新闻事件中抽离出来，放回到 DeFi 基础设施与链上行为测量的框架中理解，为后续安全研究和检测系统设计提供了一个清晰起点。

\section*{引用信息}
原始论文：Dabao Wang, Siwei Wu, Ziling Lin, Lei Wu, Xingliang Yuan, Yajin Zhou, Haoyu Wang, and Kui Ren. \textit{Towards A First Step to Understand Flash Loan and Its Applications in DeFi Ecosystem}. In Proceedings of the Ninth International Workshop on Security in Blockchain and Cloud Computing (SBC '21), 2021.

\end{document}
